\section{Experimental Protocol}

\subsection{Local Evaluation Strategy}

The repository provides two local validation pathways:
\begin{enumerate}
  \item Baseline mode (\texttt{src/infer.py}) samples validation bundles with \texttt{infer.val\_ratio} and reports \texttt{hit@K}/\texttt{recall@K}.
  \item Training mode (\texttt{src/models/retrieval\_openclip.py}) evaluates \texttt{recall@K} after each epoch, where $K=\texttt{params.recall\_k}$.
\end{enumerate}

Products are indexed once per validation pass, bundles are encoded, and top-K nearest neighbors are retrieved by cosine similarity.

\subsection{Output Artifacts}

Main artifacts generated by the training/inference workflow:
\begin{itemize}
  \item \texttt{metrics.jsonl}: epoch-level train loss and validation recall.
  \item \texttt{best.pt}: best checkpoint by recall metric.
  \item \texttt{test\_submission.csv}: final prediction file for competition upload.
  \item \texttt{val\_metrics.json}: baseline run summary and local metrics.
\end{itemize}

\subsection{Recommended Ablations}

To measure impact of major components, the following ablations are recommended:
\begin{itemize}
  \item Image-only products vs. image+text product embeddings.
  \item Bundle full-image encoding vs. detection-guided region aggregation.
  \item No offline augmentation vs. deterministic offline augmentation.
  \item Single-GPU vs. multi-GPU, with fixed seed and equal effective batch size.
\end{itemize}
